\documentclass[11pt,a4paper]{article}
\usepackage[margin=20mm]{geometry}
\usepackage{fontspec}
\setmainfont{DejaVu Serif}
\setsansfont{DejaVu Sans}
\setmonofont{DejaVu Sans Mono}[Scale=0.85]
\usepackage{polyglossia}
\setdefaultlanguage{russian}
\usepackage{graphicx}
\usepackage{hyperref}
\usepackage{enumitem}
\usepackage{fancyvrb}
\hypersetup{hidelinks}
\setlength{\parindent}{0pt}
\setlength{\parskip}{4pt}
\usepackage{titlesec}
\titleformat{\section}{\large\bfseries}{}{0pt}{}
\titlespacing*{\section}{0pt}{10pt}{4pt}
\sloppy
\setlist{nosep,leftmargin=*}
\emergencystretch=3em
\newcommand{\shot}[2]{\begin{center}\includegraphics[width=\linewidth,height=.34\textheight,keepaspectratio]{#1}\par\small #2\end{center}}
\begin{document}
\begin{center}
{\small Университет ИТМО · Фронт-энд разработка · 2026}\par
{\LARGE Лабораторная работа № 3}\par
{\large ML-платформа на Vue}\par
Светлов Ярослав, группа J3211
\end{center}
\section*{Цель и технологии}
Перенести приложение ЛР2 на Vue 3, сохранив работу с API. Использованы Vite, Vue Router, Axios, Bootstrap, Chart.js и JSON Server с авторизацией. Сохранены темы и SVG-спрайт из домашних работ.
\section*{Соответствие требованиям}
\begin{itemize}
\item Vue Router: вход, регистрация, кабинет, поиск, эксперимент по ID, реестр моделей и неизвестный маршрут. Переходы защищены проверкой сессии.
\item Общий экземпляр Axios передаёт токен. Ответ 401 очищает сессию и возвращает на страницу входа.
\item Компоненты AppNav, ExperimentTable, ModelTable, ArtifactList, MetricsChart, RequestState и SvgIcon выделяют повторяющиеся части интерфейса.
\item Composable useAuth хранит и восстанавливает сессию. useRequest управляет загрузкой и ошибками, не применяет устаревшие ответы.
\item График создаётся при монтировании компонента и уничтожается при уходе со страницы. Проверки владельца выполняются не только интерфейсом, но и сервером.
\end{itemize}

\section*{Результаты проверки}
19.09.2026 выполнен повторный проход в Firefox 148.0.2 и Chromium 146.0.7680.153: 26 сценариев Firefox и 6 дополнительных сценариев Chromium пройдены. Проверены регистрация, повторный email, неверный пароль, восстановление и завершение сессии, отдельные и совместные фильтры, сброс и пустой результат, четыре разных эксперимента, неизвестный ID, скачивание JSON/TXT, создание версии и смена статусов с перезагрузкой. Проверены недействительный токен, недоступный API и узкие окна (Firefox 500 px, Chromium 390 px). Дополнительно проверены неизвестный маршрут, смена системной темы вместе с графиком и сохранение введённых данных формы после сбоя API. Установка, сборка и девять API-тестов проходят.
\section*{Запуск и проверка}
Требуется Node.js 22.12 или новее. Команды выполняются из папки этой работы.
\begin{verbatim}
npm ci
npm run api
\end{verbatim}
Во втором терминале из той же папки:
\begin{verbatim}
npm start
\end{verbatim}
Открыть \url{http://localhost:8086}. Автотесты: \texttt{npm test}. Сборка: \texttt{npm run build}; просмотр сборки: \texttt{npm run preview} вместо dev-сервера. API должен оставаться запущенным.
 API: \url{http://localhost:3006}. Учебные аккаунты: \texttt{admin@nexus.ai / password} и \texttt{vision@nexus.ai / cv12345}. Данные сохраняются в \texttt{db.local.json}, исходная база \texttt{db.json} не меняется.\section*{Вывод и ограничения}
Функции платформы сохранены при переходе к компонентам и маршрутам Vue. Обучение и облачное развёртывание не выполняются: статус Production имитирует деплой. Артефакты содержат демонстрационные метрики и логи, а не веса моделей. JSON Server с json-server-auth предназначен только для localhost и учебных данных; нельзя использовать реальные пароли.

\section*{Скриншоты}
\shot{checks/light.png}{Рисунок 1. Эксперимент и график в светлой теме.}
\shot{checks/dark.png}{Рисунок 2. Тот же интерфейс в тёмной теме.}

\end{document}
